\documentclass[10pt,a4paper]{amsart}
\usepackage[margin=19mm]{geometry}
\usepackage{amsmath,amssymb,mathtools}
\usepackage[scr=rsfs,cal=euler]{mathalfa}
\usepackage{xfrac}
\usepackage[unicode,hidelinks,bookmarks=false]{hyperref}
\newcommand{\ie}{i.e.\ }
\newcommand{\cf}{cf.\ }
\newcommand{\resp}{respectively\ }
\newcommand{\Subsection}{\subsection}
\providecommand{\nobreakdash}{\nobreak\hbox{-}}
\setlength{\emergencystretch}{2em}
\multlinegap0pt
\usepackage{amssymb}
\usepackage{xcolor}
\newtheorem{theorem}{Theorem}
\title{The Dirichlet problem with entire data for non-hyperbolic quadratic hypersurfaces}
\author{J. M. Aldaz, H. Render}
\date{Source: arXiv:2404.16735v2 (CC BY 4.0)}
\begin{document}
\maketitle

\noindent\emph{Source and licence.} The Dirichlet problem with entire data for non-hyperbolic quadratic hypersurfaces. Version arXiv:2404.16735v2 (CC BY 4.0), distributed by arXiv under the Creative Commons Attribution 4.0 International licence, http://creativecommons.org/licenses/by/4.0/, as recorded in the arXiv licence metadata for this exact version and shown on the abstract page https://arxiv.org/abs/2404.16735v2. Full licence notice: \texttt{licenses/arxiv-2404.16735-CC-BY-4.0.txt}.\par\smallskip

\noindent\textit{Editorial scope.} Counted unit: the article's theorem arbdimBB (source0.tex lines 418--431). The article's complete original proof of it is delivered with it (source0.tex lines 433--682, one \texttt{proof} environment of 250 lines). No mathematics is written, repaired or re-derived: the statement and the proof are the article's own lines, in the article's own notation, and no retained line is substituted or re-flowed.  Retained uncounted as the source-local context the proof needs: lines 199--205; lines 211--218; lines 225--233; lines 339--346; lines 351--394.  Nothing was omitted from any retained span, so the package carries no omission record.  The retained lines cite 1 works, whose entries are transcribed verbatim from the article's own bibliography source in the e-print and delivered with short editorial labels recorded in the item's reference mapping.  No retained line contains a figure environment, an image-inclusion command or drawing code, so no image is delivered and the package declares no asset dependency.  Editorial formatting only, in addition: an AMS \texttt{amsart} preamble that restates the article's own declarations and macros used by the retained lines, namely the environments \texttt{theorem} on separate counters, together with the standard packages those lines need.  The retained environments print in this document as Theorem 1 in order of appearance, whereas the article prints the same items with the numbering of its own full document.  The complete native source of the article as distributed in the arXiv e-print for this exact version is bundled unchanged as \texttt{sources/158159/source0.tex}.
\begin{equation}
x^{2}P_{n}^{\left( \alpha ,\alpha \right) }\left( x\right) =\widetilde{a}%
_{n}\left( \alpha \right) P_{n+2}^{\left( \alpha ,\alpha \right) }\left(
x\right) +\widetilde{b}_{n}\left( \alpha \right) P_{n}^{\left( \alpha
,\alpha \right) }\left( x\right) +\widetilde{g}_{n}\left( \alpha \right)
P_{n-2}^{\left( \alpha ,\alpha \right) }\left( x\right),  \label{eqrec2}
\end{equation}%

\begin{equation}
 y \ r_{u} \left( y \right) =
 \widetilde{a}_{u}\left( \alpha \right) r_{u+1} \left(
	y \right) +\widetilde{b}_{u}\left( \alpha \right)  
r_{u} \left( y  \right)  +\widetilde{g}_{u}\left( \alpha \right)
 r_{u - 1} \left(
	 y  \right).  \label{eqrec3}
\end{equation}

\begin{theorem} \label{Elbert}
Let $x_{2n,1}\left( \alpha \right) $ be the first positive zero of $
P_{2n}^{\left( \alpha ,\alpha \right) }\left( x\right) $, where $n\geq 3$ and $\alpha \geq -1/2.$ Then 
\begin{equation*}
x_{2n,1}\left( \alpha \right) \geq \frac{\pi }{4\sqrt{\left( \alpha +\frac{1
}{2}+n\right) \left( n+2\right) }}.
\end{equation*}

\end{theorem}

\begin{equation}
Y_{k,\left( s,l\right) }^{d}\left( x\right) 
:=
Y_{s,l}^{d-1}\left(
x_{1},\dots,x_{d-1}\right) \cdot \left\vert x\right\vert ^{k-s}P_{k-s}^{\left(
s+\left( d-3\right) /2, s+\left( d-3\right) /2\right) }\left( \frac{x_{d}}{%
\left\vert x\right\vert }\right)   \label{eqintFF},
\end{equation}

 Using the abreviation $\alpha _{s}:=s+\left( d-3\right) /2$, for $k\geq 1$ identity (\ref{eqintFF}) becomes 
\begin{equation*}
Y_{k,\left( s,l\right) }^{d}\left( x\right) 
=
Y_{s,l}^{d-1}\left(
x_{1},\dots ,x_{d-1}\right) \cdot \left\vert x\right\vert
^{k-s}P_{k-s}^{\left( \alpha _{s},\alpha _{s}\right) }\left( \frac{x_{d}}{
\left\vert x\right\vert }\right).
\end{equation*}
Note  that the same factor $Y_{s,l}^{d-1}\left(
x_{1},\dots,x_{d-1}\right)$ appears in the preceding expression for different values of $k$, when $\left\vert x\right\vert =1$. 
The recurrence relation (\ref{eqrec2}) together with (\ref{eqintFF}) imply that for $k\ge 2$,
\begin{equation*}
x_{d}^{2}Y_{k,\left( s,l\right) }^{d}\left( x\right) =\widetilde{a}
_{k-s}\left( \alpha _{s}\right) Y_{k+2,\left( s,l\right) }^{d}\left(
x\right) +\widetilde{b}_{k-s}\left( \alpha _{s}\right) Y_{k,\left(
s,l\right) }^{d}\left( x\right) +\widetilde{g}_{k-s}\left( \alpha
_{s}\right) Y_{k-2,\left( s,l\right) }^{d}\left( x\right).
\end{equation*}
\color{red}  \color{black}
  Define the normalized spherical
harmonic by
\begin{equation*}
\widetilde{Y}_{k,\left( s,l\right) }^{d}=d_{k,\left( s,l\right) }\cdot
Y_{k,\left( s,l\right) }^{d}\text{ where }d_{k,\left( s,l\right)
}:=\left\Vert Y_{k,\left( s,l\right) }^{d}\right\Vert _{\mathbb{S}
^{d-1}}^{-1}.
\end{equation*}
It follows that 
\begin{eqnarray}
x_{d}^{2}\widetilde{Y}_{k,\left( s,l\right) }^{d}\left( x\right) &
=
&
\widetilde{a}_{k-s}\left( \alpha _{s}\right) \frac{d_{k,\left( s,l\right) }}{
d_{k+2,\left( s,l\right) }} 
\widetilde{Y}_{k+2,\left( s,l\right) }^{d}\left(
x\right) +\widetilde{b}_{k-s}\left( \alpha _{s}\right) \widetilde{Y}
_{k,\left( s,l\right) }^{d}\left( x\right)  \label{eqrell} \\
&&
+ \  \widetilde{g}_{k-s}\left( \alpha _{s}\right) 
\frac{d_{k,\left( s,l\right) }}{d_{k-2,\left( s,l\right) }}
\widetilde{Y}_{k-2,\left( s,l\right)}^{d}\left( x\right)  
\label{eqrell2}.
\end{eqnarray}

\begin{theorem} \label{arbdimBB}
	For all homogeneous polynomials $
	f_{m}$ of degree $m$ in $d$ variables, 
	and each $j = 1, \dots , d$, we have
	\begin{equation*}
		\left\langle x_{j}^{2}f_{m},f_{m}\right\rangle _{L^{2}\left( \mathbb{S}%
			^{d-1}\right) }
		\geq
		\frac{\pi ^{2}}{4\left( m+ 2 d + 1 \right)^{2}}
		\left
		\langle
		f_{m},f_{m}\right\rangle _{L^{2}\left( \mathbb{S}^{d-1}\right) }.
	\end{equation*}
\end{theorem}

\begin{proof} Without loss of generality we set $j = d$.
  By the preceding Lemma it suffices
to show that for all even indices we can take
$$C_{2m} := \frac{\pi^{2}}{4\left( 2m+ 2 d\right) ^{2}}.
$$


 By the Gauss
decomposition (cf. Theorem 5.5 in \cite{ABR92} or Theorem 5.7 in the 2001 edition) there exist homogeneous
harmonic polynomials $h_{2k}$ of degree $2k$ for $k=0,\dots ,m$, such that 
\begin{equation}
f_{2m} (x) =\sum_{k=0}^{m}  \left\vert x\right\vert ^{2m-2k} h_{2k}(x)\text{.}
\label{eqAlmansi1}
\end{equation}
 Each harmonic polynomial $h_{2k}$
of degree $2k$ for $k=0,\dots ,m$ is a linear combination of the orthonormal
spherical harmonics $\widetilde{Y}_{2k,\left( s,l\right) }^{d}$ with $
s=0,1,2,\dots ,2k$ and $l=1,\dots ,a_{s}^{(d-1)}$. Thus, we can write 
\begin{equation*}
h_{2k}=\sum_{s=0}^{2k}\sum_{l=1}^{a_{s}^{\left( d-1\right) }}c_{k,s,l}^{d}
\widetilde{Y}_{2k,\left( s,l\right) }^{d}
\end{equation*}
for suitable complex coefficients $c_{k,s,l}^{d}.$ \color{red} \color{black}

Recall that $\lfloor x \rfloor$ denotes the integer part of $x$. By reordering it is easy
to see that 
\begin{equation*}
f_{2m} (x) 
=
\sum_{k=0}^{m} \left\vert x\right\vert
^{2m-2k} h_{2k} (x)
=
\sum_{s=0}^{2m}\sum_{l=1}^{a_{s}^{\left( d-1\right) }}\sum_{k=\left\lfloor
\frac{s+1}{2}\right\rfloor }^{m}c_{k,s,l}^d\left\vert x\right\vert ^{2m-2k}%
\widetilde{Y}_{2k,\left( s,l\right) }^{d}  (x).
\end{equation*}%
By orthogonality 
\begin{equation}
\left\langle f_{2m},f_{2m}\right\rangle _{L^2(\mathbb{S}^{d-1})}=\sum_{s=0}^{2m}%
\sum_{l=1}^{a_{s}^{\left( d-1\right) }}\sum_{k=\left\lfloor \frac{s+1}{2}\right\rfloor
}^{m}\left\vert c_{k,s,l}^d\right\vert ^{2}  \label{eqF1}
\end{equation}
(since we are integrating over the unit sphere).
Now $\left\langle x_{d}^{2}f_{2m},f_{2m}\right\rangle _{L^2(\mathbb{S}^{d-1})}$
is equal to 
\begin{equation*}
\sum_{s_{1}=0}^{2m}\sum_{l_{1}=1}^{a_{s_1}^{\left( d-1\right) }}\sum_{k_{1}= 
\left\lfloor \frac{s_1+1}{2}\right\rfloor }^{m}\sum_{s_{2}=0}^{2m}\sum_{l_{2}=1}^{a_{s_2}^{
\left( d-1\right) }}\sum_{k_{2}=\left\lfloor \frac{s_2+1}{2}\right\rfloor
}^{m}c_{k_{1},s_{1},l_{1}}^d\overline{c_{k_{2},s_{2},l_{2}}^d}\left\langle
x_{d}^{2}\widetilde{Y}^d_{2k_{1},\left( s_{1},l_{1}\right) },\widetilde{Y}^d
_{2k_{2},\left( s_{2},l_{2}\right) }\right\rangle _{L^2(\mathbb{S}^{d-1})},
\end{equation*}
so by orthogonality (using the relation (\ref{eqrell}-\ref{eqrell2}) for $
x_d^2 \widetilde{Y}^d_{2k,\left( s,l\right) }$)
we obtain    
\begin{equation}
\left\langle x_{d}^{2}f_{2m},f_{2m}\right\rangle _{L^2(\mathbb{S}^{d-1})}=\sum_{s=0}^{2m}\sum_{l=1}^{a_{s}^{\left( d-1\right) }}\sum_{k_{1}=
\left\lfloor \frac{s+1}{2}\right\rfloor }^{m}\sum_{k_{2}=\left\lfloor \frac{s+1}{2}\right\rfloor
}^{m}c^d_{k_{1},s,l}\overline{c^d_{k_{2},s,l}}A_{k_{1},k_{2}}\left( s,l\right)
\label{eqF2}
\end{equation}%
where 
\begin{equation*}
A_{k_{1},k_{2}}\left( s,l\right) :=\left\langle x_{d}^{2}\widetilde{Y}^d
_{2k_{1},\left( s,l\right) },\widetilde{Y}^d_{2k_{2},\left( s,l\right)
}\right\rangle_{L^2(\mathbb{S}^{d-1})}.
\end{equation*}
In view of (\ref{eqF1}) and (\ref{eqF2}), it suffices to show that for each $
s=0,\dots , 2m$ and each $l=1,\dots ,a_{s}^{\left( d-1\right) }$,
\begin{equation}
\sum_{k_{1}=\left\lfloor \frac{s+1}{2}\right\rfloor }^{m}\sum_{k_{2}=\left\lfloor \frac{s+1}{2}%
\right\rfloor }^{m}c^d_{k_{1},s,l}\overline{c^d_{k_{2},s,l}}A_{k_{1},k_{2 }}\left(
s,l\right) 
\geq
 \frac{\pi ^{2}}{4\left( 2m + 2d  \right) ^{2}}\sum_{k=\left\lfloor 
\frac{s+1}{2}\right\rfloor }^{m}\left\vert c^d_{k,s,l}\right\vert ^{2}.  \label{eqF3}
\end{equation}
 The integer $s$ is either even or odd,
so we can write $s=2\sigma +\delta $ with $\delta \in \left\{ 0,1\right\}.$ 
Then 
\begin{equation*}
\left\lfloor \frac{s+1}{2}\right\rfloor =\left\lfloor \frac{2\sigma +1+\delta }{2}\right\rfloor
=\sigma +\delta .
\end{equation*}%
Let us define the symmetric $\left( m-\sigma -\delta +1\right) \times \left(
m-\sigma -\delta +1\right) $-matrix 
\begin{equation*}
A_{m}\left( s,l\right) :=\left( A_{k_{1},k_{2}}\left( s,l\right) \right)
_{k_{1},k_{2}=\sigma +\delta ,\dots , m}.
\end{equation*}
Writing $c^t = \left(c^d_{\left\lfloor \frac{s +1}{2}\right\rfloor, s,l}, \dots,  c^d_{m,s,l}\right)$, it is clear that the left hand side of (\ref{eqF3}) equals 
$$
c^t \ A_m(s,l) \ \overline{c},
$$
so it is enough to show that  the  smallest eigenvalue $\lambda _{m, s,l}^{\ast }$  of
 $A_m \left( s,l\right) $
satisfies the inequality 
$$
\lambda _{m, s,l}^{\ast }
\ge
\pi ^{2}/4\left( 2m+2d\right) ^{2}.
$$
Since $A_{m}\left( 2\sigma +1,l\right) $ is the submatrix of $A_{m}\left(
2\sigma ,l\right)$ obtained by deleting the first row and the first column, its lowest eigenvalue is at least as large as $\lambda _{m, 2\sigma,l}^{\ast }$. Thus, it suffices to prove (\ref{eqF3}) for $
A_{m}\left( 2\sigma ,l\right) $, or equivalently, for  $\delta =0.$



We now use the recurrence relation (\ref{eqrell} - \ref{eqrell2}) (with $2 k$ 
instead of  $k$)  
\begin{eqnarray*}
x_{d}^{2}\widetilde{Y}_{2k,\left( s,l\right) }^{d}\left( x\right) &=&%
\widetilde{a}_{2k-s}\left( \alpha _{s}\right) \frac{d_{2k,\left( s,l\right) }%
}{d_{2k+2,\left( s,l\right) }}\widetilde{Y}_{2k+2,\left( s,l\right)
}^{d}\left( x\right) +\widetilde{b}_{2k-s}\left( \alpha _{s}\right) 
\widetilde{Y}_{2k,\left( s,l\right) }^{d}\left( x\right) \\
&&+\widetilde{g}_{2k-s}\left( \alpha _{s}\right) \frac{d_{2k,\left(
s,l\right) }}{d_{2k-2,\left( s,l\right) }}\widetilde{Y}_{2k-2,\left(
s,l\right) }^{d}\left( x\right) .
\end{eqnarray*}

From the orthonormality relations we see that 
\begin{equation*}
A_{k,k}\left( s,l\right) =\left\langle x_{d}^{2}\widetilde{Y}^d_{2k,\left(
s,l\right) },\widetilde{Y}^d_{2k,\left( s,l\right) }\right\rangle_{L^2(\mathbb{S}^{d-1})}
=\widetilde{b}_{2k-s}\left( \alpha _{s}\right), 
\end{equation*}
\begin{equation*}
A_{k,k+1}\left( s,l\right) =\left\langle x_{d}^{2}\widetilde{Y}^d_{2k,\left(
s,l\right) },\widetilde{Y}^d_{2k+2,\left( s,l\right) }\right\rangle _{L^2(\mathbb{S}^{d-1})}
=
\widetilde{a}_{2k-s}\left( \alpha _{s}\right) \frac{d_{2k,\left(
s,l\right) }}{d_{2k+2,\left( s,l\right) }},
\end{equation*}
and 
\begin{equation*}
A_{k,k-1}\left( s,l\right) =\left\langle x_{d}^{2}\widetilde{Y}^d_{2k,\left(
s,l\right) },\widetilde{Y}^d_{2k-2,\left( s,l\right) }\right\rangle _{L^2(\mathbb{S}^{d-1})}
=
\widetilde{g}_{2k-s}\left( \alpha _{s}\right) \frac{d_{2k,\left(
s,l\right) }}{d_{2k-2,\left( s,l\right) }},
\end{equation*}
for $k=\sigma ,\dots , m$ and $s=2\sigma$, with all the other entries being equal to zero.  It follows that $
A_{m}\left( s,l\right) $ is an $\left( m-\sigma +1\right) \times \left(
m-\sigma +1\right) -$matrix of the form 
\begin{equation*}
\left( 
\begin{array}{cccc}
\widetilde{b}_{0} & \widetilde{a}_{0}\frac{d_{2\sigma ,\left( s,l\right) }}{%
d_{2\sigma +2,\left( s,l\right) }} &  &  \\ 
\widetilde{g}_{2}\frac{d_{2\sigma +2,\left( s,l\right) }}{d_{2\sigma ,\left(
s,l\right) }} & \widetilde{b}_{2} & \ddots &  \\ 
& \ddots & \ddots & \widetilde{a}_{2m - 2-2\sigma }\frac{d_{2m-2,\left( s,l\right) }}{d_{2m,\left( s,l\right)
}} \\ 
&  & \widetilde{g}_{2m-2\sigma }\frac{
d_{2m,\left( s,l\right) }}{d_{2m-2,\left( s,l\right) }} & \widetilde{b}_{2m-2\sigma }
\end{array}%
\right),
\end{equation*}
  We recall some known facts about tridiagonal $n\times n$-matrices of the form 
\begin{equation*}
J_{n}=\left( 
\begin{array}{cccc}
\beta _{0} & \alpha _{0} &  &  \\ 
\gamma _{1} & \beta _{1} & \ddots &  \\ 
& \ddots & \ddots & \alpha _{n-2} \\ 
&  & \gamma _{n-1} & \beta _{n-1}%
\end{array}
\right). 
\end{equation*}%
where all the entries $\alpha _{k}$ and $\gamma _{k}$ are non-zero. In order to compute
the characteristic polynomial $\det \left( J_{n}-\lambda I_n \right)$
 of $J_{n}$,  we define  a sequence of polynomials $p_{0}\left( x\right) :=1$, 
 $
 p_{1}\left( x\right) :=\left( x-\beta _{0}\right) /\alpha _{0}$, 
 and for $0<k<n$,   we use the recurrence relation  
 \begin{equation}
 	\label{genrec}
 	xp_{k}\left( x\right) =\gamma _{k}p_{k-1}\left( x\right) +\beta
 	_{k}p_{k}\left( x\right) +\alpha _{k}p_{k+1}\left( x\right).
 \end{equation}
 Then an inductive argument for $n\ge 1$, together with the  expansion of the determinant  along the last column, yields that $p_{n}\left( x\right) $ is a constant multiple
of the characteristic polynomial of $J_{n}$, 
and more specifically,  that
\begin{equation*}
\det \left( J_{n}-\lambda I_n   \right) =\left( -1\right) ^{n}\alpha
_{0}\dots \alpha _{n-2}p_{n}\left( \lambda \right) .
\end{equation*}%
Thus the eigenvalues of $J_{n}$ are the zeros of the polynomials $
p_{n}\left( x\right)$ defined by (\ref{genrec}).  Next we consider the special form $
A_{m}\left( s,l\right) $ of  $J_n$ that appears in our argument, with an additional simplification.  Replacing $p_{k}$ with $
q_{k}:=d_{k}p_{k}$ (where $d_{k}\neq 0)$ in
(\ref{genrec}) we obtain the recurrence relation 
\begin{equation*}
xq_{k}\left( x\right) =\alpha _{k}\frac{d_{k}}{d_{k+1}}q_{k+1}\left(
x\right) +\beta _{k}q_{k}\left( x\right) +\gamma _{k}\frac{d_{k}}{d_{k-1}}
q_{k-1}\left( x\right),
\end{equation*}
with the initial conditions transformed to $q_{0}\left( x\right) =d_{0}$
and $q_{1}\left( x\right) =d_{1}\left( x-\beta _{0}\right) /\alpha _{0}.$

It follows that the matrix $A_{m}\left( s,l\right) $, for $s=2\sigma $, has the
same eigenvalues as the matrix
\begin{equation*}
\left( 
\begin{array}{cccc}
\widetilde{b}_{0} & \widetilde{a}_{0} &  &  \\ 
\widetilde{g}_{2} & \widetilde{b}_{2} & \ddots  &  \\ 
& \ddots  & \ddots  & \widetilde{a}_{2m- 2 - 2\sigma } \\ 
&  & \widetilde{g}_{2m-2\sigma } & \widetilde{b}_{2m-2\sigma }
\end{array}
\right). 
\end{equation*}%
Next we replace the even indices $2n$ appearing
in the preceding matrix with $u = 2n$, 
and we recognize it as the Jacobi matrix
whose eigenvalues are computed via the polynomials appearing in the recurrence relation (\ref{eqrec3}). 
Denote  by  $y_u$ the first positive zero  of $r_u (y) =
P_{2n}^{\left( \alpha ,\alpha \right) } \left( \sqrt{y}\right),$ 
with $n =  m-\sigma +1$ and $\alpha
=s+\left( d-3\right) /2.$  Then  the smallest eigenvalue $\lambda _{m, s,l}^{\ast }$  of $
A_{m}\left( s,l\right) $ is $y_u$.  
Note that $\sigma $ ranges from $
0$ to $m$, so $n =m-\sigma +1$ ranges from $1$ to $m+1.$ Furthermore,  
\begin{equation*}
\alpha =s+\left( d-3\right) /2=2\left( m+1-n\right) +\left( d-3\right) /2.
\end{equation*}
By Theorem \ref{Elbert},
\begin{equation*}
\left( \alpha +\frac{1}{2}+n\right) \left( n+2\right) y_u\geq \frac{\pi ^{2}}{16}.
\end{equation*}
For $n=m-\sigma +1$ we estimate     
\begin{eqnarray*}
\left( \alpha +\frac{1}{2}+n \right) \left( n+2\right)  &\leq
&\max_{n=1,\dots , m+1}\left( 2\left( m+1-n\right) +\left( d-3\right) /2 + 1/2 +    n \right)
\cdot \left( n+2\right)  \\
&\leq & \left( m + d \right)^{2},
\end{eqnarray*}
using the fact that $d \ge 2$, which in particular entails that for every $\sigma \ge 0$ we have $3 \sigma - \sigma^2 - (d \sigma)/2 \le 1$. 

It follows that  
\begin{equation*}
\lambda _{m, s,l}^{\ast } 
=
 y_{u} 
\geq 
\frac{\pi ^{2}}{16\left( m+ d\right)^{2}}.
\end{equation*}
\end{proof}

\begin{thebibliography}{99}
\bibitem[ABR92]{ABR92} S. Axler, P. Bourdon, W. Ramey, \emph{Harmonic Function Theory,} Springer, New York 1992.
\end{thebibliography}
\end{document}
